% !PRINTED 163
\weberpage{163}
%<<CONT
wurzeln, und $e$ soll der \gesperrt{Index der Invariante} $\mathit\Psi(x)$ heissen.
Ist $\varepsilon$ eine primitive $e^{\mathrm{te}}$ Einheitswurzel, so können wir
\[
\alpha=\varepsilon^a,\quad \beta=\varepsilon^b,\quad \gamma=\varepsilon^c \ldots
\]
setzen, und die Exponenten $a,b,c\ldots$ können keinen gemeinschaftlichen
 Theiler mit $e$ haben. Daraus folgt, dass man die
ganzen Zahlen $x,y,z,\ldots$ so bestimmen kann, dass
\[
ax+by+cz+\cdots \equiv 1\ (\mathrm{mod}\ e)
\]
wird (Bd. I, \S.\,118). Da nun wegen der Gruppennatur unter
den Factoren $\alpha,\beta,\gamma,\ldots$ auch die Zahl
\[
\alpha^x\,\beta^y\,\gamma^z \cdots = \varepsilon
\]
vorkommt, so folgt, dass die Gesammtheit der Factoren $\alpha,\beta,\gamma\ldots$
mit den Potenzen von $\varepsilon$:
\[
1,\ \varepsilon,\ \varepsilon^2,\ \ldots\ \varepsilon^{e-1}
\]
zusammenfallen muss.

Suchen wir in $S$ alle Substitutionen $A$, die der Bedingung
\weq{4}{\mathit\Psi[A(x)] = \mathit\Psi(x)}
genügen, zu denen gewiss die identische Substitution gehört, so
erhalten wir eine neue Gruppe $T$, die ein Theiler von $S$ ist.
Bedeutet dann $E$ eine der Bedingung
\weq{5}{\mathit\Psi[E(x)] = \varepsilon\,\mathit\Psi(x)}
genügende Substitution aus $S$, so haben alle Substitutionen $AE$
und $EA$ die gleiche Eigenschaft, und wir erhalten die Zerlegung
von $S$ in die Nebengruppen:
\weq{6}{S = T + TE + TE^2 + \cdots + TE^{e-1};}
der Index des Theilers $T$ ist also $=\varepsilon$. % [sic]
Zugleich ergiebt sich
noch, da jede Substitution $E^{-1}AE$ der Bedingung (4) genügt,
\weq{7}{E^{-1}TE = T,}
woraus hervorgeht, dass $T$ ein Normaltheiler von $S$ ist. Wir
haben damit den Satz bewiesen:

3. \gesperrt{Ist $\mathit\Psi(x)$ eine Invariante der Gruppe $S$ vom
Index $e$, so bilden alle Substitutionen von $S$,
durch die $\mathit\Psi(x)$ ungeändert bleibt, einen Normaltheiler
 $T$ von $S$ vom Index $e$.}

Für die Gruppe $T$ ist $\mathit\Psi(x)$ absolute Invariante. Die Invarianten
 vom Index 1 sind die absoluten Invarianten von $S$.
Die Gruppe $T$ heisst die \gesperrt{zur Invariante $\mathit\Psi(x)$ geh\"orige
Gruppe.}
